\documentclass{article}
\usepackage[margin=2cm]{geometry}
\usepackage{amsmath}
\usepackage{graphicx}
\usepackage{booktabs}
\usepackage[utf8]{inputenc}
\usepackage{ragged2e}
\setlength{\emergencystretch}{3em}
\begin{document}
\sloppy

allows us to conclude that this is always the case for real initial conditions \$8. Thus we
have a one parameter family of real smooth solutions, labeled by the IR parameter \$8. With
this in mind, we may choose any value of \$8 and solve the BPS equations in numerically.
In Figure , we plot the solutions obtained for the following choices of initial condition:.
\$8 = \{0.25, 0.5, 1, 1.5, 2\}. In order to get smooth solutions for u > 0, we must take n = -1.
It is straighforward to verify that the resulting solutions are completely smooth and have
the expected vanishing of e2 at the origin, implying that the spacetime smoothly pinches
off. Furthermore, e2f /e2u is seen to asymptote to a constant, which we denote by e2fe..


\subsection*{0.1UV asymptotic expansions}


As in the holographic Janus solutions in Lorentzian signature , the BPS equations may also
be used to obtain the UV asymptotic behavior of the solutions. To do so, we begin by
defining an asymptotic coordinate z = e-", where the asymptotic S boundary is reached by
coefficients in the UV expansions of the non-zero fields may now be solved for order-by-order.
using the BPS equations. One finds explicitly that all coefficients are determined in terms of.
only three independent parameters a, , and fx, in accord with the fact that there are three
independent first-order differential equations. The first few terms in the expansions are


\begin{equation}
\begin{array} { c } { \dot { f ( z ) = - \log z + f _ { k } - \left ( \frac { 1 } { 4 } e ^ { - 2 f _ { k } } + \frac { 1 } { 1 6 } \alpha ^ { 2 } \right ) ~ z ^ { 2 } + O ( z ^ { 4 } ) } } \\ { \dot { \sigma ( z ) = \frac { 3 } { 8 } \alpha ^ { 2 } ~ z ^ { 2 } + \frac { 1 } { 4 } e ^ { I _ { k } } \alpha \beta ~ z ^ { 3 } + O ( z ^ { 4 } ) } } \\ { \dot { \phi ^ { 0 } ( z}}\end{array}
\end{equation}


We have obtained the expansions up to O(z8), but we display only the first few terms here.


\subsection*{1Holographic sphere free energy.}


The goal of this section is to obtain the holographic free energy, i.e. the renormalized on-shell
action. We begin by writing the full action,.


\begin{equation}
\begin{array} { r } { S _ { \mathrm { G D } } = \int d u ~ d ^ { 5 } x ~ \sqrt { G } ~ \mathcal { L } } \\ { S _ { \mathrm { a s h, o n t e ~ A l e s u s i c a s i n s i o n t e s t ~ D e a n s i o n t ~ D e a n d e s a n s i c a l ~ D e a n d e s s ~ C a n d e s s ~ a n d e s i n s ~ a n d e s ~ a n d e s ~ a n d e s ~ a n d e s ~ I}}}\end{array}
\end{equation}


where S6p is the six-dimensional Euclidean action given in and SGh is the Gibbons-Hawking
term. The y appearing in SGh is the determinant of the induced metric on the boundary

\end{document}
